% !TEX root = ../main.tex
\chapter*{Conclusion générale}
\addcontentsline{toc}{chapter}{Conclusion générale}
\markboth{Conclusion générale}{Conclusion générale}
\justifying

Ce rapport présente le travail accompli lors de ce projet de fin d'études, réalisé au sein du service des Relations Internationales de Toulouse INP-ENSEEIHT dans le cadre du cycle ingénieur en génie informatique de l'ENICarthage. Il m'a permis de mettre en pratique l'ensemble des connaissances théoriques acquises durant ma formation, depuis l'analyse d'un besoin métier jusqu'à la livraison d'une application complète.

Dans un premier temps, j'ai mené une analyse approfondie du contexte du projet et des solutions existantes, afin de proposer une solution adaptée aux contraintes du service RI. En m'appuyant sur la méthodologie Scrum, j'ai ensuite planifié et réalisé six sprints successifs, chacun couvrant un périmètre fonctionnel cohérent, de la gestion des référentiels et des accords de mobilité jusqu'au moteur de recommandation de destinations. Chaque sprint a été suivi d'une phase de tests visant à valider le travail réalisé.

Ce projet m'a permis de renforcer mes compétences en conception logicielle, à travers la modélisation UML, la justification argumentée de choix algorithmiques et la conception de machines à états pour sécuriser le cycle de vie des entités métier. J'ai également découvert et mis en œuvre de nouvelles technologies, notamment Django Ninja et Next.js, ainsi qu'un modèle d'apprentissage supervisé pour le moteur de recommandation.

J'ai utilisé GitHub tout au long du projet, aussi bien pour la gestion du code source que pour le suivi du backlog via les Issues et les Milestones, et je me suis appuyé sur une chaîne d'intégration continue (GitHub Actions) exécutant les tests et les contrôles de qualité à chaque évolution du code.

Ce travail a permis de livrer une plateforme fonctionnelle répondant aux objectifs fixés : les cinq sprints du MVP ont été réalisés et validés, ainsi que le moteur de recommandation de destinations (US-22) du Sprint~6. Certaines pistes constituent des perspectives d'évolution pour une prochaine version, notamment l'assistant virtuel fondé sur une architecture RAG (US-23), qui n'a pas pu être mené à terme.

Sur le plan de la validation, trois points restent en deçà de l'objectif initial. Le comportement de l'API sous forte charge n'a pas pu être mesuré de manière concluante, les essais ayant été limités par l'environnement de test disponible ; le seuil de temps de réponse visé n'est donc pas confirmé. L'audit de sécurité dynamique, réalisé avec OWASP~ZAP sur l'interface et l'API, ne révèle aucune alerte de niveau élevé, mais il a été mené sans authentification et hors conditions de production. Enfin, la couverture des tests d'interface reste inégale et la disponibilité du système n'a pas fait l'objet d'une mesure dédiée. Une campagne de charge dans un environnement représentatif, un audit de sécurité approfondi et un renforcement des tests d'interface constituent les premières tâches d'une future itération.

Ces évolutions futures permettraient de renforcer davantage la plateforme et de répondre à des besoins complémentaires du service des Relations Internationales.
